\chapter{Derivation T: The Rough Path Lift for Continuum Observables}
\label{ch:derivation_t_rough_path}

\section{Context: The Continuum Observable Problem}
In the continuum limit ($\beta \to \infty$), the gauge field $A_\mu$ becomes a distribution of regularity $C^{-1/2-\epsilon}$ (like white noise in 2D, even worse in 4D). The classical line integral $W_C = \text{Tr } \mathcal{P} \exp \oint A$ is mathematically ill-defined because the product of distributions is ambiguous.
We use \textbf{Rough Path Theory} (Gubinelli/Hairer) to rigorously define the Wilson loop as a continuous functional on a space of rough paths, removing lattice regularization artifacts. This relies on \textbf{Appendix AG}.

\section{Theorem T.1 (Continuity of the Holonomy Map)}
The Wilson loop functional can be lifted to a continuous map on the space of **Rough Paths** $\mathbf{A} = (A, \mathbb{A})$, where $A$ is the field and $\mathbb{A}$ is the second-order iterated integral (Levy area).
\begin{equation}
\Phi: (A, \mathbb{A}) \mapsto \text{Holonomy}(A)
\end{equation}
is well-defined and continuous in the $p$-variation topology for $p < 3$.

\section{Proof Derivation}
\subsection{Step 1: Signature Construction}
We define the "signature" of the gauge field on lattice segments. Using Balaban’s bounds (Derivation J), we verify the probabilistic Hölder condition:
\begin{equation}
|A_{s,t}| \le C |t-s|^\alpha
\end{equation}
and for the second iterated integral (lattice plaquette):
\begin{equation}
|\mathbb{A}_{s,t}| \le C |t-s|^{2\alpha}
\end{equation}

\subsection{Step 2: Sewing Lemma}
We apply Gubinelli’s Sewing Lemma. Since the Hölder exponent $\alpha > 1/3$ (in effective dimension), we have $2\alpha + \alpha > 1$. This satisfies the analytical condition required to construct the global path integral from local approximations.

\subsection{Step 3: Renormalized Trace}
The "naive" calculation of the trace diverges due to the self-energy of the loop. We define the renormalized trace using the rough path geometric correction (counterterm for the self-intersection of Brownian motion).
\begin{equation}
\langle W_{ren} \rangle = \lim_{a \to 0} Z(a)^{-1} \langle W_{latt}(a) \rangle
\end{equation}

\subsection{Step 4: Limit Existence}
This proves that the continuum observable $W(C)$ exists as a random variable in the physical Hilbert space, independent of the lattice spacing. This effectively solves the problem of defining "local" observables in a theory of distributions.
